\chapter{Identitätsbeziehung zwischen DSTT und Weber-Kräften}

\section{Einleitung}

In Kapitel 8 wird die Behauptung aufgestellt, dass die Zustandsvariable der 
Dynamischen Schwere-Trägheits-Theorie (DSTT) und die Terme der 
Weber-Gravitation (WG) über eine Identitätsbeziehung verknüpft sind:

\[
\frac{|\vec{F}_{\text{WG}}^{\text{tan}}|}{|\vec{F}_{\text{WG}}^{\text{rad}}| + |\vec{F}_{\text{WG}}^{\text{tan}}|} \stackrel{?}{=} \beta_{\text{DSTT}}(t)
\]

Dieser Anhang zeigt, dass diese Beziehung **nicht momentan** gilt. 
Die korrekte Identität ist eine Beziehung zwischen den zeitlich gemittelten 
Leistungen und der Drift von \(\beta_{\text{DSTT}}\).

\section{Definitionen}

\subsection{DSTT-Zustandsvariable}

Nach Kapitel 8.2.3 ist der Zustandsparameter der DSTT definiert als das 
Verhältnis von tangentialer zu gesamter kinetischer Energie:

\[
\beta_{\text{DSTT}}(t) = \frac{E_{\text{tan}}(t)}{E_{\text{kin}}(t)} = 
\frac{\frac{1}{2} m r^2 \dot{\phi}^2}{\frac{1}{2} m \dot{r}^2 + \frac{1}{2} m r^2 \dot{\phi}^2}
\tag{G.1}
\]

\subsection{Weber-Gravitation für Massen}

Die Weber-Kraft für massive Körper (Kapitel 8.4.3 mit \(\beta_{\text{WG}} = 0,5\)):

\[
\vec{F}_{\text{WG}} = -\frac{GMm}{r^2} \left\{ 
\left[1 - \frac{\dot{r}^2}{c^2} + \frac{1}{2} \frac{r\ddot{r}}{c^2} \right] \hat{r} + 
\frac{2(\dot{r} \cdot \vec{v})}{c^2} \vec{v} \right\}
\tag{G.2}
\]

\subsection{Radiale und tangentiale Komponenten}

Die radiale Kraftkomponente ist:

\[
F_{\text{rad}} = -\frac{GMm}{r^2} \left[ 1 - \frac{\dot{r}^2}{c^2} + \frac{1}{2} \frac{r\ddot{r}}{c^2} \right]
\tag{G.3}
\]

Die tangentiale Kraft (Vektor) ist:

\[
\vec{F}_{\text{tan}} = -\frac{GMm}{r^2} \cdot \frac{2\dot{r}^2}{c^2} \vec{v}
\tag{G.4}
\]

Der Betrag der tangentialen Kraft ist:

\[
|\vec{F}_{\text{tan}}| = \frac{2GMm}{r^2} \cdot \frac{\dot{r}^2}{c^2} \cdot v, \quad 
v = \sqrt{\dot{r}^2 + r^2\dot{\phi}^2}
\tag{G.5}
\]

\section{Warum die momentane Identität scheitert}

Ein einfaches Gegenbeispiel genügt: Für eine reine Kreisbahn (\(\dot{r} = 0\)) 
ist \(|\vec{F}_{\text{tan}}| = 0\), aber \(\beta_{\text{DSTT}} = 1\). Also:

\[
\frac{|\vec{F}_{\text{tan}}|}{|\vec{F}_{\text{rad}}| + |\vec{F}_{\text{tan}}|} = 0 \neq 1 = \beta_{\text{DSTT}}
\tag{G.6}
\]

Die momentane Identität ist daher **falsch**.

\section{Die korrekte Identität über Energiebilanz}

\subsection{Leistung der tangentialen Kraft}

Die momentane Leistung der tangentialen Kraft ist:

\[
P_{\text{tan}} = \vec{F}_{\text{tan}} \cdot \vec{v} = \frac{2GMm}{r^2} \cdot \frac{\dot{r}^2}{c^2} \cdot v^2
\tag{G.7}
\]

\subsection{Zeitliche Änderung der tangentialen Energie}

Die tangentiale kinetische Energie ist \(E_{\text{tan}} = \frac{1}{2} m r^2 \dot{\phi}^2\). 
Ihre zeitliche Ableitung ist:

\[
\frac{dE_{\text{tan}}}{dt} = m r \dot{\phi} ( \dot{r} \dot{\phi} + r \ddot{\phi} ) = 
\vec{F}_{\text{tan}} \cdot \vec{v} + \text{(Beitrag der radialen Kraft)}
\tag{G.8}
\]

Nach Mittelung über eine Umlaufperiode \(T\) verschwinden die oszillatorischen 
Anteile, und es bleibt:

\[
\left\langle \frac{dE_{\text{tan}}}{dt} \right\rangle_T = \langle P_{\text{tan}} \rangle_T
\tag{G.9}
\]

\subsection{Beziehung zu \(\dot{\beta}_{\text{DSTT}}\)}

Mit \(E_{\text{tan}} = \beta_{\text{DSTT}} E_{\text{kin}}\) folgt:

\[
\frac{dE_{\text{tan}}}{dt} = \dot{\beta}_{\text{DSTT}} E_{\text{kin}} + \beta_{\text{DSTT}} \dot{E}_{\text{kin}}
\tag{G.10}
\]

Im zeitlichen Mittel über eine Periode ist \(\langle \dot{E}_{\text{kin}} \rangle_T = 0\), 
da die kinetische Energie nach jedem Umlauf zum Ausgangswert zurückkehrt 
(keine säkulare Änderung auf dieser Zeitskala). Somit:

\[
\langle \dot{\beta}_{\text{DSTT}} \rangle_T = \frac{\langle P_{\text{tan}} \rangle_T}{\langle E_{\text{kin}} \rangle_T}
\tag{G.11}
\]

Dies ist die korrekte Identität: Die Drift von \(\beta_{\text{DSTT}}\) ist 
direkt mit der mittleren Leistung der tangentialen Weber-Kraft verknüpft.

\subsection{Numerische Abschätzung}

Für eine Kepler-Ellipse mit Exzentrizität \(e\) gilt:

\[
\langle \dot{r}^2 \rangle_T = \frac{GM}{a} \cdot \frac{e^2}{2}, \quad
\langle v^2 \rangle_T = \frac{GM}{a}
\]

Damit folgt aus (G.11):

\[
\langle \dot{\beta}_{\text{DSTT}} \rangle_T \sim \frac{GM}{c^2 a} \cdot e^2 \cdot \omega
\tag{G.12}
\]

mit \(\omega = \sqrt{GM/a^3}\). Dies ist konsistent mit der Kopplungsgleichung 
aus Kapitel 8.6.

\section{Spezialfall: Elektrodynamik (\(\beta_{\text{WED}} = 2\))}

Für die Weber-Elektrodynamik mit \(\beta_{\text{WED}} = 2\) verschwindet 
der tangentiale Term vollständig:

\[
\vec{F}_{\text{tan}}^{\text{(WED)}} = 0 \quad \Rightarrow \quad P_{\text{tan}} = 0
\tag{G.13}
\]

Damit ist \(\langle \dot{\beta}_{\text{DSTT}} \rangle_T = 0\) – die 
Elektrodynamik ändert die Energieaufteilung nicht. Sie ist reversibel.

\section{Zusammenfassung der korrigierten Beziehungen}

\begin{center}
\begin{tabular}{lcc}
\hline
 & Gravitation (WG) & Elektrodynamik (WED) \\
\hline
\(\beta_{\text{Kraft}}\) & \(0,5\) & \(2\) \\
\(\vec{F}_{\text{tan}}\) & \(\propto \dot{r}^2 \vec{v} / c^2\) & \(0\) \\
\(P_{\text{tan}}\) & \(> 0\) (außer \(\dot{r}=0\)) & \(0\) \\
\(\langle \dot{\beta}_{\text{DSTT}} \rangle_T\) & \(> 0\) (irreversibel) & \(= 0\) (reversibel) \\
\hline
\end{tabular}
\end{center}

\section{Konsequenzen für die Haupttext-Behauptungen}

\begin{enumerate}
    \item Die in Kapitel 8.5 behauptete momentane Identität 
          \(\frac{|\vec{F}_{\text{tan}}|}{|\vec{F}_{\text{rad}}| + |\vec{F}_{\text{tan}}|} \equiv \beta_{\text{DSTT}}(t)\) 
          ist **falsch** (Gegenbeispiel: Kreisbahn).
    \item Die korrekte Beziehung ist (G.11): \(\langle \dot{\beta}_{\text{DSTT}} \rangle_T = \langle P_{\text{tan}} \rangle_T / \langle E_{\text{kin}} \rangle_T\).
    \item Der qualitative Schluss (\(\beta_{\text{WG}} = 0,5 \Rightarrow \dot{\beta} > 0\)) bleibt korrekt.
    \item Die Elektrodynamik ist reversibel (\(\dot{\beta} = 0\)).
\end{enumerate}

\section{Ausblick}

Die hier präsentierte korrigierte Identität ist konsistent mit der 
Energiebilanz und der Kopplungsgleichung aus Kapitel 8.6. Eine mögliche 
zukünftige Verfeinerung wäre die Herleitung einer exakten differentiellen 
Beziehung zwischen \(\vec{F}_{\text{tan}}\) und \(\dot{\beta}_{\text{DSTT}}\) 
ohne Mittelung, die jedoch aufgrund der oszillatorischen Natur der Kräfte 
nicht trivial ist.